\begin{theorem}[Algorithm 3.1 has a globally maximum initial series]
Let \(G\) be a finite loopless multigraph satisfying the tablet's
edge-criticality condition.  Suppose that
\(\chi'(G)=k+1\) and \(\Delta(G)+1\leq k\).  Then there exist an edge
identity \(e_0\), a partial \(k\)-edge-coloring \(\varphi\), a tree-sequence
\(T\), and a certified Tashkinov series \(\mathcal T\) with zero
noninitial rounds, such that \((T,e_0,\varphi)\) is globally maximum and:
\[
 \mathcal T.T_0=T|y_1,\qquad
 \mathcal T.T_1=T,\qquad
 \mathcal T.\varphi_0=\varphi;
\]
\[
 \mathcal T.S_0=\mathcal T.F_0=\varnothing,
 \quad \mathcal T.\Theta_0=\mathrm{none},
 \quad \mathcal T.v_0=\mathcal T.u_0=
 \mathcal T.\delta_0=\mathcal T.\gamma_0=\mathrm{none}.
\]
Here \(T|y_1\) is the exact length-one prefix of \(T\).  Moreover, the
terminal slot satisfies
\[
 |V(\mathcal T.T_1)|=t(G),\qquad |V(\mathcal T.T_1)|\geq3,
 \qquad |V(\mathcal T.T_1)|\text{ is odd},
\]
and \(V(\mathcal T.T_1)\) is both elementary and closed with respect to
\(\varphi\).
\end{theorem}
\begin{proof}
By \noderef{GlobalMaximumTripleExists}, choose an edge identity \(e_0\),
a partial \(k\)-edge-coloring \(\varphi\), and a tree-sequence \(T\) such
that \((T,e_0,\varphi)\) is globally maximum.  In particular, by the
definition in \noderef{GlobalMaximumTriple}, \((G,e_0,\varphi)\) is a
\(k\)-triple and \(T\) is a Tashkinov tree for this triple.

Let \(P=T|y_1\) be the exact length-one prefix supplied by
\noderef{TreeSequencePrefix}.  The theorem
\noderef{TashkinovTreeFromLengthOnePrefix} gives a Tashkinov augmentation
sequence under \(\varphi\) from \(P\) to \(T\).  On the other hand,
\noderef{ScheideGlobalMaximumOrder} gives
\[
 t(G)=|V(T)|,\quad |V(T)|\geq3,\quad |V(T)|\text{ odd},
\]
and says that \(V(T)\) is elementary and closed with respect to
\(\varphi\).  The augmentation sequence together with this last closedness
clause is exactly a Tashkinov closure of \(P\) under \(\varphi\), by
\noderef{TashkinovClosure}.

Define zero-round series data as follows.  In the two tree slots put \(P\)
at index \(0\) and \(T\) at index \(1\).  In the unique coloring slot put
\(\varphi\).  In the unique stored-color and connecting-edge slots put the
empty sets, and put the value none in each of the unique extension-kind,
chosen-vertex, outside-vertex, delta-color, and gamma-color slots.  These
are all the fields required by \noderef{TashkinovSeriesData}, and the
displayed slot equations hold directly by construction.

It remains to certify the data.  There is only one certificate index,
namely index \(0\).  The initial constructor of
\noderef{TashkinovSeriesStep} applies there: its triple premise is supplied
by the globally maximum triple; the seed \(P\) has length one by its exact
prefix construction; the coloring, empty-set, and absent-witness premises
are the defining equations above; and its closure premise is the closure
of \(P\) by \(T\) just proved.  Thus the data and this sole certificate form
a zero-round Tashkinov series by \noderef{TashkinovSeries}.

Its terminal tree slot is definitionally \(T\).  Substitution into the
five conclusions of \noderef{ScheideGlobalMaximumOrder} yields respectively
the asserted global-order equality, lower bound, oddness, elementarity, and
closedness of the terminal slot.  This proves every claimed property of
\(\mathcal T\).
\end{proof}
